\section{Variation and principal values}
\label{sec:consequences}

The uniform-rectifiability conclusion also allows the forward estimates
of Mas and Tolsa to be applied under the individual hard-truncation
bound \eqref{eq:riesz-bound}. For a scalar or Euclidean-vector family
$F=(F_\varepsilon)_{\varepsilon>0}$ and $\rho>2$, write
\[
 V_\rho(F)(x)=
 \sup_{\substack{m\ge1\\ \varepsilon_0>\cdots>\varepsilon_m>0}}
 \left(\sum_{j=0}^{m-1}
 |F_{\varepsilon_{j+1}}(x)-F_{\varepsilon_j}(x)|^\rho
 \right)^{1/\rho}.
\]
Thus the supremum is over all finite decreasing sequences of positive
scales, with the Euclidean norm for vector values. Let
$K:\R^d\setminus\{0\}\to\R$ be an odd $C^2$ convolution kernel satisfying,
for some finite $C_K>0$,
\begin{equation}\label{eq:variation-kernel}
 K(-z)=-K(z),\qquad
 |\partial^\beta K(z)|\le C_K|z|^{-n-|\beta|}
 \quad (|\beta|\le2,\ z\ne0),
\end{equation}
and set
\[
 T^{K,\mu}_\varepsilon f(x)
 =\int_{|x-y|>\varepsilon}K(x-y)f(y)\,d\mu(y).
\]
For the smooth family, fix one nondecreasing
$\varphi_{\R}\in C^2([0,\infty);[0,\infty))$ such that, for some finite
$C_\varphi>0$,
\begin{equation}\label{eq:variation-cutoff}
 \mathbf1_{[4,\infty)}\le\varphi_{\R}
 \le\mathbf1_{[1/4,\infty)},\qquad
 \mathbf1_{[1/3,3]}\le C_\varphi|\varphi_{\R}'|.
\end{equation}
Put $\varphi_\varepsilon(z)=\varphi_{\R}(|z|^2/\varepsilon^2)$ and
\[
 T^{K,\mu}_{\varphi,\varepsilon}f(x)
 =\int K(x-y)\varphi_\varepsilon(x-y)f(y)\,d\mu(y),
\]
where the cutoff kernel is defined to be zero on the diagonal.

\begin{corollary}[Consequences of Mas--Tolsa]\label{cor:variation}
Let $d\ge4$ and $2\le n\le d-2$ be integers, and let $\mu$ be an
$n$-Ahlfors--David regular Radon measure satisfying
\eqref{eq:riesz-bound}. Fix $\rho>2$ and a real odd $C^2$ convolution
kernel $K$ satisfying \eqref{eq:variation-kernel}. For every real
$f\in L^2(\mu)$,
\begin{equation}\label{eq:hard-variation-consequence}
 \bigl\|V_\rho((T^{K,\mu}_\varepsilon f)_{\varepsilon>0})\bigr\|_{L^2(\mu)}
 \le C_2\|f\|_{L^2(\mu)}.
\end{equation}
The principal value
\[
 T^{K,\mu}f(x)=\lim_{\varepsilon\downarrow0}T^{K,\mu}_\varepsilon f(x)
\]
exists finitely for $\mu$-almost every $x$, belongs to $L^2(\mu)$, and
\[
 \|T^{K,\mu}_\varepsilon f-T^{K,\mu}f\|_{L^2(\mu)}
 \longrightarrow0\qquad(\varepsilon\downarrow0).
\]
The vector Riesz family $R_{\mu,\varepsilon}f$ satisfies the analogous
hard-variation bound and has an almost-everywhere principal-value limit,
with strong convergence in $L^2(\mu;\R^d)$. The exceptional null set
may depend on $K$ and $f$. These are strong-convergence assertions on
each input; no operator-norm convergence is asserted.

For the fixed profile \eqref{eq:variation-cutoff}, every real
$f\in L^p(\mu)$ with $1<p<\infty$ satisfies
\begin{equation}\label{eq:smooth-variation-consequence}
 \bigl\|V_\rho((T^{K,\mu}_{\varphi,\varepsilon}f)_{\varepsilon>0})
 \bigr\|_{L^p(\mu)}
 \le C_p\|f\|_{L^p(\mu)}.
\end{equation}
For every real $f\in L^1(\mu)$ and $\lambda>0$,
\[
 \mu\left\{x:
 V_\rho((T^{K,\mu}_{\varphi,\varepsilon}f)_{\varepsilon>0})(x)
 >\lambda\right\}
 \le \frac{C_1}{\lambda}\|f\|_{L^1(\mu)}.
\]
The constants may depend on the dimensions, the fixed $\rho$, the kernel
bounds, the uniform-rectifiability data, and, for the smooth estimates,
the chosen profile and $p$ when relevant.
\end{corollary}

\begin{proof}
If $D_E=0$, the support is empty or a singleton. All the truncated
operators above vanish $\mu$-almost everywhere, so the assertions are
immediate. Suppose $D_E>0$. Theorem~\ref{thm:main} gives uniform
rectifiability, with constants $\theta,M$ and admissible radii when the
support is bounded. For unbounded support, the forward clauses
$(a)\Rightarrow(c)$ and $(a)\Rightarrow(b)$ of
Mas--Tolsa's Theorem~2.3 give, respectively, the hard $L^2$ and smooth
$L^p$/weak-$(1,1)$ estimates; their Definition~2.1 is precisely the
fixed cutoff above, and their Theorem~1.3 includes the vector Riesz
family \cite{MasTolsa2014Variation}. Their range $1\le n<d$ contains
the present one. We reduce bounded support to this unbounded setting.

Suppose $D=D_E<\infty$, and fix $a\in E$. Since $n<d$, choose an
affine $n$-plane $P$ with $\dist(a,P)=4D$, and set
\[
 \sigma=\mu+\H^n\!\restriction P.
\]
For $x\in E$, the diameter bound gives
$3D\le\dist(x,P)\le5D$. Thus $E$ and $P$ are disjoint and
$\supp\sigma=E\cup P$ is unbounded. Lemma~\ref{lem:global-growth}
and the Euclidean plane-volume bound give
\[
 \sigma(B(z,r))\le(9^dC_{\rm AD}+\omega_n)r^n
 \qquad(z\in\R^d,\ r>0),
\]
where $\omega_n$ is the volume of the unit ball in $\R^n$.

We verify lower growth and Lipschitz pieces at every support center
and radius. For centers on $P$, the plane supplies both properties at
every radius. For $x\in E$ and $r\le D$, use the corresponding
properties of $\mu$. If $D<r<12D$, the ball $B(x,D/2)$ lies in
$B(x,r)$ and has mass at least
\[
 C_{\rm AD}^{-1}(D/2)^n\ge C_{\rm AD}^{-1}24^{-n}r^n.
\]
The uniform-rectifiability map at scale $D/2$ captures at least
$\theta(D/2)^n\ge\theta24^{-n}r^n$ in the same ball. Composing
that map with $u\mapsto Du/(2r)$ gives a map on $B_{\R^n}(0,r)$
with the same image and Lipschitz constant at most $M$.
If $r\ge12D$, the disk in $P$ centered at $\pi_Px$ with radius
$r/2$ lies in $B(x,r)$, since every point $z$ of that disk satisfies
\[
 |z-x|\le |z-\pi_Px|+\dist(x,P)
 <r/2+5D\le11r/12.
\]
This disk has $\sigma$-mass at least $\omega_n2^{-n}r^n$ and is
the image of $B_{\R^n}(0,r)$ under a $1/2$-Lipschitz affine map.
These estimates prove that $\sigma$ is globally $n$-AD regular and
uniformly $n$-rectifiable, with constants controlled by the original
dimensions, AD data, $\theta$, and $M$.

For any input $f$ in one of the stated $L^p(\mu)$ spaces, extend it
by zero on $P$ to obtain $\widetilde f$. Disjointness gives
$\|\widetilde f\|_{L^p(\sigma)}=\|f\|_{L^p(\mu)}$.
For every positive truncation and every $x\in E$,
\[
 T^{K,\sigma}_\varepsilon\widetilde f(x)
   =T^{K,\mu}_\varepsilon f(x),\qquad
 T^{K,\sigma}_{\varphi,\varepsilon}\widetilde f(x)
   =T^{K,\mu}_{\varphi,\varepsilon}f(x).
\]
The vector Riesz truncations have the same identity. Apply the forward
estimates to $\sigma$ and restrict the output to $E$. Since
$\sigma|_E=\mu$, the identities transfer both the strong bounds and
the weak-$(1,1)$ bound to $\mu$, with the required constant dependence.

In either diameter regime, for a real $f\in L^2(\mu)$,
Lemma~\ref{lem:global-growth} and the size bound in
\eqref{eq:variation-kernel} give
\[
 \int_{|x-y|>\varepsilon}|K(x-y)|^2\,d\mu(y)
 \le C\varepsilon^{-n}.
\]
Consequently the displayed hard integrals converge absolutely for
every $x$ and tend to zero as $\varepsilon\to\infty$.
Approximate $f$ in $L^2$ by compactly supported $L^1\cap L^2$ inputs
$f_j$. The same kernel estimate gives pointwise convergence of each
positive truncation. Taking the supremum over finite sequences gives
\[
 V_\rho((T^{K,\mu}_\varepsilon f)_{\varepsilon>0})
 \le\liminf_{j\to\infty}
 V_\rho((T^{K,\mu}_\varepsilon f_j)_{\varepsilon>0}),
\]
so Fatou's lemma transfers the cited variation bound to the actual
integrals for $f$. For the smooth formulas, the analogous exterior
$L^{p'}(\mu)$ kernel estimate, with $p'=p/(p-1)>1$, and H\"older's
inequality give the same identification for $L^p$ inputs. For $L^1$
inputs the cutoff kernel is bounded at each fixed positive scale.

Let $V=V_\rho((T^{K,\mu}_\varepsilon f)_{\varepsilon>0})$. By
\eqref{eq:hard-variation-consequence}, $V$ is finite almost everywhere.
Finite variation forces the truncations to be Cauchy as
$\varepsilon\downarrow0$: otherwise a decreasing sequence with
infinitely many jumps bounded below would make $V$ infinite.
This proves existence of the finite principal value. The definition
gives $|T^{K,\mu}_\varepsilon f-T^{K,\mu}_\delta f|\le V$ for every
$\varepsilon,\delta>0$. Sending $\delta\downarrow0$ gives
$|T^{K,\mu}_\varepsilon f-T^{K,\mu}f|\le V$, while sending
$\delta\to\infty$ gives $|T^{K,\mu}_\varepsilon f|\le V$.
Thus the limit belongs to $L^2(\mu)$, and dominated convergence with
the majorant $V^2\in L^1(\mu)$ proves strong $L^2$ convergence.
The same argument with Euclidean norms applies to the vector Riesz
family.
\end{proof}
